\documentclass{article}
\usepackage{graphicx} % Required for inserting images

\title{Sutton and Barton Exercises Chapter 11}
\author{danieltocco0 }
\date{March 2026}

\documentclass{article}
\usepackage{amsmath}
\begin{document}


\maketitle

\section{Chapter 11}

\textit{Exercise 11.1} Convert the equation of n-step on-policy TD (7.9) to semi-gradient form.
Give accompanying definitions of the return for both the episodic and continuing cases.
\\
\\

We are pretty much just rewriting things in terms of function approximation...

\[
\mathbf{w}_{t+n} \leftarrow \mathbf{w}_{t+n-1} + \alpha \, \rho_{t:t+n-1}
\bigl[G_{t:t+n} - \hat{v}(S_t, \mathbf{w}_{t+n-1})\bigr]
\nabla \hat{v}(S_t, \mathbf{w}_{t+n-1})
\]

where the return is defined as:

\textit{Episodic:}
\[
G_{t:t+n} \doteq R_{t+1} + \gamma R_{t+2} + \cdots + \gamma^{n-1} R_{t+n}
+ \gamma^n \hat{v}(S_{t+n}, \mathbf{w}_{t+n-1})
\]

\textit{Continuing:}
\[
G_{t:t+n} \doteq R_{t+1} - \bar{R}_t + R_{t+2} - \bar{R}_{t+1}
+ \cdots + R_{t+n} - \bar{R}_{t+n-1}
+ \hat{v}(S_{t+n}, \mathbf{w}_{t+n-1})
\]
\\\\
\textbf{Exercise 11.2} \quad The semi-gradient form of $n$-step $Q(\sigma)$ is:
\\
\\

\[
\mathbf{w}_{t+n} \leftarrow \mathbf{w}_{t+n-1} + \alpha
\bigl[G_{t:t+n} - \hat{q}(S_t, A_t, \mathbf{w}_{t+n-1})\bigr]
\nabla \hat{q}(S_t, A_t, \mathbf{w}_{t+n-1})
\]

where the $Q(\sigma)$ return is:

\textit{Episodic:}
\[
G_{t:t+n} \doteq \hat{q}(S_t, A_t, \mathbf{w}) + \sum_{k=t}^{t+n-1} \delta_k
\prod_{i=t+1}^{k} \Bigl(\gamma \sigma_i \rho_i + \gamma(1-\sigma_i)\pi(A_i|S_i)\Bigr)
\]

\textit{Continuing:}
\[
G_{t:t+n} \doteq \hat{q}(S_t, A_t, \mathbf{w}) + \sum_{k=t}^{t+n-1} \delta_k
\prod_{i=t+1}^{k} \Bigl(\gamma \sigma_i \rho_i + \gamma(1-\sigma_i)\pi(A_i|S_i)\Bigr)
\]

where in the continuing case the TD error uses the differential form:
\[
\delta_k = R_{k+1} - \bar{R} + \hat{q}(S_{k+1}, A_{k+1}, \mathbf{w})
- \hat{q}(S_k, A_k, \mathbf{w})
\]
\\\\
\textbf{Exercise 11.4} Prove (11.24).:
\\
\\


\textbf{Derivation:}

\text{Start with the definition of } $\overline{RE}$:

\[
\overline{RE}(\mathbf{w}) = \mathbb{E}\left[(G_t - \hat{v}(S_t, \mathbf{w}))^2\right]
\]

\text{Add and subtract } $v_\pi(S_t)$ \text{ inside the bracket:}

\[
= \mathbb{E}\left[(G_t - v_\pi(S_t) + v_\pi(S_t) - \hat{v}(S_t, \mathbf{w}))^2\right]
\]

\text{Let } $a = G_t - v_\pi(S_t)$ \text{ and } $b = v_\pi(S_t) - \hat{v}(S_t, \mathbf{w})$\text{, expand } $(a+b)^2$:

\[
= \mathbb{E}\left[(G_t - v_\pi(S_t))^2\right] 
+ 2\mathbb{E}\left[(G_t - v_\pi(S_t))(v_\pi(S_t) - \hat{v}(S_t,\mathbf{w}))\right] 
+ \mathbb{E}\left[(v_\pi(S_t) - \hat{v}(S_t,\mathbf{w}))^2\right]
\]

\text{The cross term vanishes because } $\mathbb{E}[G_t - v_\pi(S_t) \mid S_t] = 0$ 
\text{ by definition of } $v_\pi(S_t) = \mathbb{E}[G_t \mid S_t]$:

\[
\mathbb{E}\left[(G_t - v_\pi(S_t))(v_\pi(S_t) - \hat{v}(S_t,\mathbf{w}))\right]
= \mathbb{E}\left[\mathbb{E}[G_t - v_\pi(S_t)\mid S_t] \cdot (v_\pi(S_t) - \hat{v}(S_t,\mathbf{w}))\right] = 0
\]

\text{What remains:}

\[
\overline{RE}(\mathbf{w}) 
= \mathbb{E}\left[(v_\pi(S_t) - \hat{v}(S_t,\mathbf{w}))^2\right] 
+ \mathbb{E}\left[(G_t - v_\pi(S_t))^2\right]
= \overline{VE}(\mathbf{w}) + \mathbb{E}\left[(G_t - v_\pi(S_t))^2\right]
\]

\text{The second term is the variance of } $G_t$ \text{ around the true value } $v_\pi(S_t)$
\text{ and does not depend on } $\mathbf{w}$\text{. Therefore } $\overline{RE}$ \text{ and } $\overline{VE}$ 
\text{ share the same optimal parameter vector } $\mathbf{w}^*$.
\end{document}
